
choisirons des indéterminées différentes pour chacun des facteurs du produit de \(\Ghat\) par lui-même, et nous indiquerons l'indéterminée en exposant lorsque ce sera nécessaire : \(\Ghat^t\), \(\Ghat^u\), etc. Pour \(A\in\operatorname{Ob}(\mathcal C_0)\), on a donc \(\Ghat^u\circ\Ghat^t(A)=1+\{1+A[[t]]^+\}[[u]]^+\). Pour définir les opérations \(\lambda^i\) sur \(1+K[[t]]^+\), nous allons donc nous donner un morphisme de foncteurs : \(\Ghat\longrightarrow\Ghat^u\circ\Ghat^t\).
Pour cela, il faut d'abord définir une filtration du foncteur \(\Ghat^u\circ\Ghat^t\) qui permette le recollement de 1.6.

On remarque d'abord que \(F^k\Ghat\) est un idéal pour la loi \(\circ\)~; pour cela on est ramené par fonctorialité au cas \(A=\mathbb Z[X_i,Y]_{i\geq1}\), et au produit \(\left(1+\sum_{i\geq k}X_i t^i\right)\circ(1+Yt)\), et cela résulte de 2.3.2. On prend alors pour \(F^n(\Ghat^u\circ\Ghat^t)\) l'ensemble des séries de la forme \(1+\psi_1u+\ldots+\psi_i u^i+\ldots\), \(\psi_i\in F^{[n/i]}\Ghat\), \([n/i]\) désignant la partie entière de \(n/i\).

L'élément définissant le morphisme \(\lambda_u\) est \(\underline{1}+\{1+Xt\}u\)~; le lecteur vérifiera à titre d'exercice que les conditions de 1.6 sont bien remplies (en fait, la filtration est choisie ad hoc).

Ici encore, les opérations \(\lambda^i\) se calculent par substitution à partir de \(\lambda_u(1+\ldots+X_nt^n+\ldots)\), et sont donc déterminées par des polynômes universels. On notera \(Q_{i,j}(X_1,\ldots,X_{ij})\) le coefficient de \(t^j\) dans \(\lambda^i(1+\ldots+X_nt^n+\ldots)\)~; c'est un polynôme isobare de poids \(ij\) par rapport aux variables \(X_1,\ldots,X_{ij}\) (vérification analogue à celle de 1.15). Il se calcule par fonctions symétriques à partir de l'élément de définition :
\begin{equation}\tag{2.3.3}
  \lambda_u(1+Xt)=\underline{1}+\{1+Xt\}u.
\end{equation}

Donnons à titre d'exemple la valeur de \(Q_{2,2}\) :
